\documentclass[a4paper,12pt]{article}
\usepackage{amsmath,amssymb}

\begin{document}
	
	\section*{Rationals, Irrationals, and Completeness}
	
	The real numbers are the union of rationals and irrationals:
	\[
	\mathbb{R} = \mathbb{Q} \,\cup\, (\mathbb{R} \setminus \mathbb{Q}),
	\]
	where $\mathbb{Q} = \left\{ \tfrac{p}{q} : p,q \in \mathbb{Z}, \; q \neq 0 \right\}$ and  
	$\mathbb{R} \setminus \mathbb{Q}$ are the \emph{irrationals}.
	
	\bigskip
	
	\textbf{Intuition:}  
	Inside $\mathbb{Q}$, certain numbers appear as ``holes.''  
	For example,
	\[
	A = \{q \in \mathbb{Q} : q^2 < 2\}
	\]
	is bounded above, but in $\mathbb{Q}$ it has no supremum, since $\sqrt{2} \notin \mathbb{Q}$.
	
	\bigskip
	
	\textbf{Completeness:}  
	The Axiom of Completeness ensures that in $\mathbb{R}$ every nonempty bounded set has a least upper bound.  
	Thus,
	\[
	\sup \{x \in \mathbb{R} : x^2 < 2\} = \sqrt{2} \in \mathbb{R}.
	\]
	
	\bigskip
	
	\textbf{Conclusion:}  
	Irrationals are not merely ``holes we approach'' with rationals; by completeness they are genuine members of $\mathbb{R}$.  
	This is what makes $\mathbb{R}$ continuous, unlike $\mathbb{Q}$.
	
		
		\section*{Axiom of Completeness}
		
		\textbf{Formal statement (quantifiers):}
		
		If $A \subseteq \mathbb{R}$, $A \neq \varnothing$, and $A$ is bounded above, then
		\[
		\exists\, s \in \mathbb{R} \;\; \text{such that} \;\;
		\begin{cases}
			\forall a \in A, & a \leq s, \\[6pt]
			\forall t < s, & \exists a \in A \;\; (t < a).
		\end{cases}
		\]
		Here $s = \sup A$ is the \emph{least upper bound} of $A$.
		
		\bigskip
		
		\textbf{Examples:}
		
		\begin{itemize}
			\item \textbf{Rational supremum:}
			\[
			B = \{x \in \mathbb{R} : x < 3\}, 
			\qquad \sup B = 3 \in \mathbb{Q}.
			\]
			In interval notation:
			\[
			B = (-\infty, 3).
			\]
			
			\item \textbf{Irrational supremum:}
			\[
			A = \{x \in \mathbb{R} : x^2 < 2\}, 
			\qquad \sup A = \sqrt{2} \notin \mathbb{Q}.
			\]
			In interval notation:
			\[
			A = (-\sqrt{2}, \sqrt{2}).
			\]
		\end{itemize}
		
		\bigskip
		
		\textbf{Equivalence of Notations:}
		
		\[
		(a,b) = \{x \in \mathbb{R} : a < x < b\}, \quad
		[a,b] = \{x \in \mathbb{R} : a \leq x \leq b\}.
		\]
		
		\[
		[a,b) = \{x \in \mathbb{R} : a \leq x < b\}, \quad
		(a,b] = \{x \in \mathbb{R} : a < x \leq b\}.
		\]
		
% Irrational supremum
\[
A = \{q \in \mathbb{Q} : q^2 < 2\}, 
\qquad 
\boxed{\sup A = \sqrt{2} \in \mathbb{R}}
\]

% Rational supremum
\[
B = \{x \in \mathbb{R} : x < 3\}, 
\qquad 
\boxed{\sup B = 3 \in \mathbb{Q}}
\]
The Axiom of Completeness guarantees that every bounded set has a supremum in $\mathbb{R}$, whether it is rational or irrational. 

\section*{Why We Use Sets in Analysis}

\textbf{Reason 1: Generalization.}  
Working with \emph{sets} allows us to handle whole families of numbers at once.
For example, instead of a single number, we may consider the interval
\[
(0,1) = \{x \in \mathbb{R} : 0 < x < 1\}.
\]
This collection has many elements, and we can meaningfully speak about its
infimum and supremum:
\[
\inf(0,1) = 0, \qquad \sup(0,1) = 1.
\]

\bigskip

\textbf{Reason 2: Bounds and Limits.}  
The language of sets is essential for stating results like the Axiom of Completeness:
if $A \subseteq \mathbb{R}$ is nonempty and bounded above, then $\sup A$ exists in $\mathbb{R}$.
This applies no matter how large or small $A$ is.

\bigskip

\textbf{Reason 3: Single Elements as Sets.}  
Even an individual number can be represented as a set, called a \emph{singleton}:
\[
\{5\} = \{x \in \mathbb{R} : x = 5\}.
\]
In this case,
\[
\inf\{5\} = \sup\{5\} = 5.
\]
Thus a single element is just a very small set.

\bigskip

\textbf{Your View, Formalized.}  
- If we want to study a whole collection, we place it inside a set and examine its limits and bounds.  
- If we want to study a single element, we may also represent it as a set.  
Both fit into the same formal framework:
\[
\text{Set = Collection of numbers (possibly just one).}
\]

	
	
\end{document}
